\documentclass[11pt]{article}
\usepackage[margin=0.9in]{geometry}
\usepackage{mathptmx}
\usepackage[hidelinks]{hyperref}
\pagestyle{empty}
\setlength{\parindent}{0pt}
\setlength{\parskip}{0.42\baselineskip}

\begin{document}

Avery Karlin\\
Department of Applied Physics and Applied Mathematics, Columbia University\\
Department of Computer Science, University of Colorado Boulder\\
\href{mailto:ak5232@columbia.edu}{ak5232@columbia.edu}

June 10, 2026

To the Editors of \emph{Nonlinear Dynamics},

Please consider the enclosed manuscript, ``Chimera Collapse Ages: Topology-Dependent Finite-Size Scaling in Mean-Field and Ring Oscillator Systems,'' for publication as an original research article in \emph{Nonlinear Dynamics}.

The manuscript reports that the first-passage criterion standard throughout the chimera-lifetime literature measures the wrong observable. Extending trajectories past their first criterion-satisfying crossing shows that up to 98\% of sustained near-synchrony excursions are self-healing grazes the chimera survives. An absorption-grade label roughly doubles true lifetimes and inverts the stability picture of the deployed real-time system that motivated the study: a parameter-free application of the reduced two-population flow, validated against its source's published bifurcations, places the operating corner beyond a homoclinic bifurcation --- the operated ``chimera'' is a transient on a destroyed limit cycle's ghost, while the corner discarded as unreliable holds the system's only attractor.

Four results demonstrate the significance and novelty of the work. First, the graze-versus-absorption correction applies wherever a chimera grazes the synchronization boundary, as the deployed two-population system does, and carries a direct engineering consequence for collapse-detector design; the ring, which does not graze, provides a clean control in which first-passage and absorption-grade labels coincide. Second, the corrected statistics reveal collapse as an aging process --- a hazard rising along each trajectory, expressed as a breath-synchronized ratchet of the envelope --- that survives stratification on the collective initial condition. Third, a matched, pre-registered 7,500-run campaign on the nonlocally coupled ring, whose existence boundary we measure independently at $\beta_c = 0.137$, shows the aging is topology-independent, with Weibull shape saturating to a finite $k_\infty \approx 1.35$--$1.65 > 1$ in the thermodynamic limit, while the finite-size scaling is not: near-boundary ring lifetimes fall with $N$ where mean-field lifetimes are flat over a $16\times$ range. Fourth, the same parameter-free reduction that explains the geometry of death underpredicts finite-$N$ lifetimes by a fixed $3.2\times$, with additive collective noise, multiplicative breath-locked noise, and the off-manifold Watanabe--Strogatz constants each ruled out by direct test --- finite-size fluctuations prolong this chimera's life, a sharply constrained open problem.

The work sits squarely in territory \emph{Nonlinear Dynamics} has actively developed, from traveling chimeras in coupled pendula to pinning control of chimera states, and its bifurcation backbone (saddle-node, Hopf, homoclinic, Takens--Bogdanov) and pre-registered, data-driven survival methodology match the journal's stated scope.

The manuscript is original, is not under consideration elsewhere, and has a single author. Data and code are openly available and archived with a Zenodo DOI; every figure regenerates from committed configurations. Use of generative-AI tools (Anthropic Claude and Claude Code) in implementing analysis code and in drafting is disclosed in the Acknowledgments; all results, derivations, and citations were independently verified by the author, who takes sole accountability for the work.

Thank you for your consideration.

Sincerely,

Avery Karlin

\end{document}
